% LR law: paragraph and figure caption (draft).
% Every number below is in results/paper/lr_law_numbers.json, written by
% scripts/paper_figures/fig_lr_law.py from:
%   analyses/scaling/results/lr_law_fit_cliff.json   (the law, 38 points)
%   analyses/scaling/results/lr_law_fit.json         (strict variant, 20 points)
%   results/lr_curves/lr_search_ep1_lim39959808_d352_se6_*.json (production sweep)
% Figure: results/paper/figures/fig_lr_law.{pdf,png}

\paragraph{Learning-rate law.}
Along our IsoFLOP curves the model size $N$ and the number of training examples
$D$ change together, and the ratio $D^*/N^*$ falls with compute, so a learning
rate tuned as a function of $N$ alone would conflate the two. We therefore tuned
the learning rate at every point of every IsoFLOP curve: a grid of seven rates
per $(N, D)$ cell spaced by a factor of 1.85, extended with further rates wherever
the optimum fell at an edge, with AdamW ($\beta_2 = 0.999$), batch size 256 and a
cosine schedule held fixed. Each cell's optimum is the argmin of an
Akima interpolation of validation $\dE$ against $\log\eta$, following
Porian et al.~\cite{porian2024}. Of 42 cells, 38 bracket their optimum: 20 with
stable rates on both sides, and 18 closed from above by a rate that diverged.
The remaining 4 have their best rate at the edge of the grid, which bounds the
optimum without locating it, and are excluded. Within one IsoFLOP curve $N$ and
$D$ are almost perfectly anti-correlated, so it is the six curves at different
compute budgets that separate the two exponents. A plane in log space describes
the 38 optima (Fig.~\ref{fig:lr_law}):
\begin{equation}
  \eta^*(N, D) = 4.2\times10^{5}\, N^{-0.69}\, D^{-0.72}, \qquad r^2 = 0.83,
\end{equation}
over $N$ from $8.1\times10^{4}$ to $1.8\times10^{7}$ parameters and $D$ from
$3.5\times10^{5}$ to $3.6\times10^{7}$ examples. The root-mean-square residual is
0.26 decades, a factor of 1.8, about one step of the tuning grid; a law in $N$
alone explains much less of the variation ($r^2 = 0.50$). The exponents are
stable: refitting with any one compute budget left out gives $N$ exponents from
$-0.70$ to $-0.66$ and $D$ exponents from $-0.74$ to $-0.63$, and the law fitted
without the highest budget predicts that budget's five optima with a median
ratio of measured to predicted of 1.00 (worst 5.1). Requiring a stable rate on
both sides of every optimum leaves 20 points and gives $N^{-0.49} D^{-0.74}$
($r^2 = 0.74$): the $D$ exponent is unchanged and the $N$ exponent shallower.
We use the 38-point law because it spans every model size in the sweep.
As a held-out check, we swept eight rates at the production configuration
($N = 1.13\times10^{7}$, $D = 4.0\times10^{7}$, 1.1 times the largest $D$ in the
fit), which the fit never saw. The law predicts $\eta^* = 1.9\times10^{-5}$; the
sweep's lowest validation cross-entropy is at $1.94\times10^{-5}$ and its lowest
greedy $\dE$ at $9.7\times10^{-6}$, one grid step (a factor of 2) below the
prediction (single seed).
% OPEN, decide before submission: the production checkpoint was trained at
% 3.1e-4, the highest rate that did not diverge, 16x above the law. On the same
% sweep it has greedy dE 12.12 vs 9.12 at 9.7e-6, and val CE 6.205 vs 6.013 at
% 1.94e-5. The paper needs one sentence on why (or a retrain at the law's rate).

\begin{figure}
  \centering
  \includegraphics[width=\textwidth]{fig_lr_law}
  \caption{\textbf{The tuned learning rate lies close to one plane in
  $(\log N, \log D, \log\eta)$.} Each marker is the $\dE$-optimal learning rate
  of one IsoFLOP point, at its own model size $N$ and number of training
  examples $D$. Circles: optimum bracketed by stable rates on both sides.
  Triangles: optimum closed from above by a rate that diverged. Open circles:
  best rate at the edge of the grid, a bound, excluded from the fit. The
  translucent plane is the fitted law,
  $\eta^* = 4.2\times10^{5} N^{-0.69} D^{-0.72}$ ($r^2 = 0.83$, 38 points);
  vertical lines are residuals to it. Star: the $\dE$-optimal rate of a separate
  sweep at the production configuration, not used in the fit.
  \textbf{a}, oblique view: the rate falls along both $N$ and $D$.
  \textbf{b}, view along the plane: scatter about it.}
  \label{fig:lr_law}
\end{figure}
